%%%%%%%%%%%%%%%%%
% CV targeted for TerX / OneSide Technologies - Ingénieur Développement Moyens d'Essais
%%%%%%%%%%%%%%%%%

\documentclass[8pt,a4paper,ragged2e,withhyper]{../_shared/altacv}

\geometry{left=1.25cm,right=1.25cm,top=.5cm,bottom=1.cm,columnsep=1.2cm}

\usepackage{paracol}
\usepackage{fontawesome5}
\usepackage{hyperref}

\iftutex 
  \setmainfont{Roboto Slab}
  \setsansfont{Lato}
  \renewcommand{\familydefault}{\sfdefault}
\else
  \usepackage[rm]{roboto}
  \usepackage[defaultsans]{lato}
  \renewcommand{\familydefault}{\sfdefault}
\fi

\definecolor{SlateGrey}{HTML}{2E2E2E}
\definecolor{LightGrey}{HTML}{666666}
\definecolor{DarkPastelRed}{HTML}{8AC0D9}
\definecolor{PastelRed}{HTML}{00B0DA}
\definecolor{GoldenEarth}{HTML}{B4AD0C}

\colorlet{name}{black}
\colorlet{tagline}{PastelRed}
\colorlet{heading}{DarkPastelRed}
\colorlet{headingrule}{GoldenEarth}
\colorlet{subheading}{PastelRed}
\colorlet{accent}{PastelRed}
\colorlet{emphasis}{SlateGrey}
\colorlet{body}{LightGrey}

\definecolor{violet}{HTML}{5A2582}
\definecolor{rouge}{HTML}{F53B3B}
\definecolor{noir}{HTML}{00232C}
\definecolor{bleu}{HTML}{00B0DA}
\definecolor{rose}{HTML}{F831A0}
\definecolor{jaune}{HTML}{FFEC00}
\definecolor{vert}{HTML}{34AD6C}

\renewcommand{\namefont}{\Huge\rmfamily\bfseries}
\renewcommand{\personalinfofont}{\footnotesize}
\renewcommand{\cvsectionfont}{\LARGE\rmfamily\bfseries}
\renewcommand{\cvsubsectionfont}{\large\bfseries}
\renewcommand{\cvItemMarker}{{\small\textbullet}}
\renewcommand{\cvRatingMarker}{\faCircle}

\input{../_shared/pubs-num.tex}
\addbibresource{../_shared/sample.bib}

\begin{document}

\name{BOILEAU Guillaume}
\tagline{Ingénieur Moyens d'Essais | Bancs de test, acquisition, LabVIEW, signal \& validation}

\photoL{2.8cm}{../_shared/moi}

\personalinfo{%
  \email{guillaume.boileaupro@gmail.com}
  \phone{+33 6 59 94 85 84}
  \location{Alpes-Maritimes, FRANCE}
  \linkedin{guillaume-boileau-phd-891b81265}
  \researchgate{Guillaume-Boileau-2}
  \orcid{0000-0002-3576-6968}
}

\makecvheader
\vspace{-0.3cm}

\AtBeginEnvironment{itemize}{\small}
\columnratio{0.64}

\begin{paracol}{2}

\cvsection{Experience}

\cvevent{\textbf{Software Engineer -- System Integration, Hardware Interfaces \& Operational Diagnostics}}{VEV Platform Services France}{2025 -- 2026}{Valbonne, France}

\textbf{Context:} Development and integration of software services for an EV charging CPMS platform connected to field equipment, operational data flows and hardware infrastructure.

\textbf{Objective:} Support reliable software integration, operational data-flow continuity, interface diagnostics and technical follow-up for services connected to charging station systems.

\begin{itemize}
  \item \textbf{Software development:} Developed and maintained web and backend services using TypeScript, Angular and Node.js in an operational industrial environment.
  \item \textbf{System integration:} Worked at the interface between software services, field equipment constraints, operational data and production expectations.
  \item \textbf{Hardware/software interfaces:} Analysed interactions between platform services, connected charging stations, data exchanges and operational behaviours.
  \item \textbf{Operational data acquisition:} Monitored and analysed field-related data flows, including FTP-based exchanges used for technical follow-up.
  \item \textbf{Diagnostics:} Investigated incidents involving software behaviour, data consistency, interface continuity and equipment-related operational effects.
  \item \textbf{Validation logic:} Checked coherence between data flows, application behaviour, equipment states and expected operational results.
  \item \textbf{Monitoring \& tooling:} Used Git, Zenhub, Confluence, OpenSearch and Sentry for issue tracking, diagnostic follow-up and technical documentation.
\end{itemize}

\divider

\cvevent{\textbf{Senior R\&D Engineer -- Test Logic, Signal Processing \& Validation Workflows}}{CNES, Laboratoire Lagrange, Observatoire de la Côte d'Azur}{2023 -- 2025}{Nice, France}

\textbf{Context:} Engineering activities for the LISA ESA/NASA space mission, involving complex software chains, simulation outputs, signal analysis, performance assessment and validation campaigns.

\textbf{Objective:} Develop, integrate and validate Python-based workflows for model comparison, signal/data analysis, consistency checks and traceable technical review.

\begin{itemize}
  \item \textbf{Validation workflows:} Structured comparison, verification and validation logic for heterogeneous simulation outputs and technical deliverables.
  \item \textbf{Signal processing:} Analysed weak, noisy and superposed signals using spectral methods, signal/noise separation logic and performance metrics.
  \item \textbf{Measurement-oriented analysis:} Worked on simulated and instrument-related data where performance depends on noise, sensitivity, acquisition quality and signal interpretation.
  \item \textbf{Automated test logic:} Defined repeatable analysis sequences, comparison rules, consistency checks and acceptance-oriented diagnostics.
  \item \textbf{Data reliability:} Ensured traceability of inputs, configurations, outputs, assumptions, metadata and validation results.
  \item \textbf{Software platform:} Developed CATpyLISA, a Python platform for model integration, data comparison, validation and technical review.
  \textcolor{DarkPastelRed}{\href{https://gitlab.oca.eu/gboileau/lisa_l3_tools}{\bf GitLab:CATpyLISA}}
  \item \textbf{Technical documentation:} Used Confluence, GitLab, Jupyter and LaTeX to document methods, validation results and technical evidence.
  \item \textbf{Coordination:} Coordinated technical contributions, supervised interns and supported collaborative validation activities in an international mission environment.
\end{itemize}

\divider

\cvevent{\textbf{Data Scientist -- Noise Diagnostics, Measurement Analysis \& Performance Validation}}{Universiteit Antwerpen}{2021 -- 2023}{Antwerp, Belgium}

\textbf{Context:} Data analysis and performance studies for precision instruments, with focus on noisy measurements, environmental disturbances and sensitivity limitations.

\textbf{Objective:} Identify, characterize and mitigate sources affecting measurement quality and system sensitivity using reproducible statistical and signal-processing workflows.

\begin{itemize}
  \item \textbf{Test data analysis:} Analysed measurement and simulation outputs to identify abnormal behaviours, noise contributions and degradation sources.
  \item \textbf{Signal \& noise diagnostics:} Developed methods for spectral analysis, noise characterization, sensitivity studies and cross-sensor correlations.
  \item \textbf{Python pipelines:} Built reproducible analysis workflows for data cleaning, transformation, modelling, diagnostics and reporting.
  \item \textbf{Performance validation:} Assessed system sensitivity, measurement stability and limitations through statistical and signal-based criteria.
  \item \textbf{Technical reporting:} Produced traceable analyses and technical summaries for international engineering and research stakeholders.
\end{itemize}

\divider

\cvevent{\textbf{Junior R\&D Engineer -- Instrumental Data, Simulation \& Signal Analysis}}{ARTEMIS, Observatoire de la Côte d'Azur}{2018 -- 2021}{Nice, France}

\textbf{Context:} R\&D activities linked to weak-signal analysis, numerical simulation, instrumental diagnostics and gravitational-wave mission preparation.

\textbf{Objective:} Develop analysis methods and simulation workflows for complex, noisy and large-scale datasets.

\begin{itemize}
  \item \textbf{Signal analysis:} Developed methods for the analysis and separation of weak, overlapping signals in noisy environments.
  \item \textbf{Simulation workflows:} Built numerical workflows for repeated analysis campaigns, uncertainty studies and performance assessment.
  \item \textbf{Instrumental diagnostics:} Investigated noise sources and limiting factors using correlation analysis and signal-processing methods.
  \item \textbf{Validation evidence:} Compared models, simulations and outputs to assess consistency, robustness and technical limitations.
\end{itemize}

\switchcolumn

\cvsection{Profile for TerX}

\textbf{Engineer at the interface between test means, software development, signal processing, acquisition-oriented data analysis and validation of complex systems}. Background developed in space, precision instrumentation and industrial software environments, with strong experience in measurement-quality analysis, automated workflows, traceability and technical diagnostics.

Relevant for a role in \textbf{test bench software, instrumentation, acquisition, automated test sequences and performance validation}: able to understand physical measurements, structure robust analysis logic, develop technical software, investigate anomalies, document results and work with multidisciplinary teams.

\cvsection{Fit for Test Means}

\begin{itemize}
  \item Strong match with \textbf{signal processing}, measurement analysis, noise characterization and performance assessment.
  \item Experience building \textbf{automated analysis workflows} and repeatable validation sequences in Python.
  \item Familiarity with \textbf{LabVIEW} from academic training; ready to reactivate and deepen it in an industrial test-bench context.
  \item Good understanding of \textbf{instrumentation logic}: acquisition, signal quality, data reliability, calibration-oriented reasoning, diagnostics and exploitation of measurements.
  \item Experience with \textbf{software/hardware interfaces} through operational systems connected to field equipment.
  \item Strong practice of \textbf{technical documentation}, result traceability, validation reports and collaborative engineering tools.
  \item Comfortable in R\&D, production-support and multidisciplinary environments involving software, data, instrumentation and system constraints.
\end{itemize}

\cvsection{Technical Skills}

\cvtag{Python}
\cvtag{LabVIEW}
\cvtag{Test logic}
\cvtag{Signal processing}
\cvtag{Data acquisition}
\cvtag{Measurement analysis}
\cvtag{Noise analysis}
\cvtag{Spectral analysis}
\cvtag{Automation}
\cvtag{Validation}
\cvtag{IVVQ}
\cvtag{V-cycle}
\cvtag{Test benches}
\cvtag{Instrumentation}
\cvtag{Hardware/software interfaces}
\cvtag{Data reliability}
\cvtag{Traceability}
\cvtag{Technical reporting}
\cvtag{Linux}
\cvtag{Git}
\cvtag{GitLab}
\cvtag{Docker}
\cvtag{Confluence}
\cvtag{Zenhub}
\cvtag{OpenSearch}
\cvtag{Sentry}
\cvtag{Jupyter}
\cvtag{SQL}
\cvtag{C/C++}
\cvtag{Matlab}
\cvtag{TypeScript}
\cvtag{Angular}
\cvtag{Node.js}

\cvsection{Methods \& Engineering}

\begin{itemize}
  \item Test sequences, test logic, integration checks, verification and validation criteria.
  \item Acquisition/data-flow monitoring, consistency checks, anomaly diagnostics and root-cause-oriented analysis.
  \item Signal/noise analysis, SNR reasoning, performance validation and robustness assessment.
  \item Technical documentation, validation evidence, reporting and traceability of assumptions/results.
  \item Agile delivery, task tracking, technical coordination and collaborative engineering.
\end{itemize}

\cvsection{Languages}

\cvskill{English}{4}
\divider

\cvskill{Spanish}{2}
\divider

\medskip

\cvsection{Education}

\cvevent{PhD in Physics}{Université Côte d'Azur}{2018 -- 2021}{Nice, France}
\divider

\cvevent{MSc in Astronomy, Astrophysics \& Space Engineering}{Observatoire de Paris / Université Paris-Saclay}{2017 -- 2018}{Paris, France}
\divider

\cvevent{Selective Graduate Program in Fundamental Physics}{Université Paris-Saclay}{2016 -- 2018}{Orsay, France}
\divider

\cvevent{BSc in Fundamental Physics}{Université Paris-Saclay}{2015 -- 2016}{Orsay, France}

\cvsection{Professional Training}

\cvevent{Supervised Machine Learning}{DeepLearning.AI -- Andrew Ng}{2020}{Coursera}
\divider

\cvevent{Recherche reproductible : principes méthodologiques pour une science transparente}{FUN MOOC -- Inria}{2025}{France Université Numérique}
\divider

\cvevent{Continuous Professional Training -- Software, Data, AI and Systems}{OpenClassrooms}{2024 -- 2026}{}

\small{Python, SQL, Machine Learning, AI, Linux, JavaScript, TypeScript, Angular, Node.js, Django, REST APIs, C/C++, Java, algorithms and software engineering fundamentals.}

\end{paracol}

\cvsection{Projects}

\cvevent{CNES / LISA -- Validation, Signal Processing \& Automated Analysis Workflows}{ESA/NASA Space Mission -- Large-scale International Collaboration}{2023 -- 2025}{Nice, France}

Contribution to software, signal-processing and validation workflows for the LISA space mission, with emphasis on traceable comparison of outputs, noisy signal analysis, performance diagnostics and consolidation of technical evidence.

\textbf{Relevance for test means:} automated validation logic, comparison sequences, signal/noise analysis, consistency checks, traceable outputs, technical documentation and multidisciplinary coordination.

\divider

\cvevent{Precision Instrumentation -- Noise Diagnostics \& Sensitivity Studies}{Universiteit Antwerpen / Gravitational-wave Instrumentation Context}{2021 -- 2023}{Antwerp, Belgium}

Development of statistical and signal-processing workflows to identify noise sources, evaluate measurement quality and support performance improvement for precision instruments.

\textbf{Relevance for instrumentation:} measurement analysis, spectral diagnostics, cross-sensor correlations, data reliability, technical reporting and performance validation.

\divider

\cvevent{Operational Industrial Platform -- Hardware-connected Software Services}{VEV Platform Services France}{2025 -- 2026}{Valbonne, France}

Software development and diagnostic support for a CPMS platform connected to EV charging infrastructure, involving operational data flows, production incidents and hardware/software interface analysis.

\textbf{Relevance for industrial testing:} interface diagnostics, data-flow monitoring, incident analysis, software reliability and operational traceability.

\end{document}
